% Generisano: uređuj beleske/sNNN.tex i naslove u manifestu.
\documentclass[11pt,a4paper]{article}
% Zajednički preambulum za dokument za čitanje; nije Beamer tema.
\usepackage[a4paper,margin=20mm,headheight=15pt]{geometry}
\usepackage{lmodern,fontspec}
\usepackage[serbian]{babel}
\setmainfont{Latin Modern Sans}
\setsansfont{Latin Modern Sans}
\setmonofont{Latin Modern Mono}
\usepackage{amsmath,amssymb,mathtools,graphicx,xcolor}
\usepackage{array,booktabs,tabularx,multirow,makecell,listings,fancyhdr}
\pagestyle{fancy}\fancyhf{}
\newcommand{\BeleskeTekuciID}{NEOZNACEN}
\fancyhead[L]{\small Beleške uz slajd \texttt{\expandafter\detokenize\expandafter{\BeleskeTekuciID}}}
\fancyfoot[R]{\small\thepage}
\setlength{\parindent}{0pt}\setlength{\parskip}{6pt}
\newwrite\BeleskeMapa
\immediate\openout\BeleskeMapa=\jobname.notes-map.tsv
\AddToHook{shipout/before}{%
  \immediate\write\BeleskeMapa{\BeleskeTekuciID\space\thepage}}
\AddToHook{enddocument/afterlastpage}{\immediate\closeout\BeleskeMapa}
\newcommand{\BeleskeID}[1]{%
  \def\BeleskeZadatiID{#1}%
  \ifx\BeleskeTekuciID\BeleskeZadatiID\else
    \PackageError{beleske}{Pogresan ID beleske: #1}{Proveri mapiranje slajda.}%
  \fi}
\newcommand{\BeleskaSlajda}[4]{%
  \clearpage\gdef\BeleskeTekuciID{#1}%
  {\Large\bfseries Slajd \textnormal{\texttt{\detokenize{#1}}}: #3\par}\medskip
  \includegraphics[page=#2,width=\linewidth]{\BeleskePrezentacija}\par\medskip
  \input{#4}}

\newcommand{\BeleskePrezentacija}{build/09_1_mos_logicka_kola.pdf}
\begin{document}
\BeleskaSlajda{09_1-s001}{1}{Logička kola sa MOS tranzistorima — 1. deo}{beleske/s001.tex}
\BeleskaSlajda{09_1-s002}{2}{Simboli MOS tranzistora}{beleske/s002.tex}
\BeleskaSlajda{09_1-s003}{3}{NMOS sa indukovanim dugačkim kanalom}{beleske/s003.tex}
\BeleskaSlajda{09_1-s004}{4}{Modulacija dužine kanala i granica modela}{beleske/s004.tex}
\BeleskaSlajda{09_1-s005}{5}{Brzina nosilaca: dugačak i kratak kanal}{beleske/s005.tex}
\BeleskaSlajda{09_1-s006}{6}{Efekat kratkog kanala}{beleske/s006.tex}
\BeleskaSlajda{09_1-s007}{7}{Parametri modela kratkog kanala}{beleske/s007.tex}
\BeleskaSlajda{09_1-s008}{8}{Dodatna aproksimacija brzine nosilaca}{beleske/s008.tex}
\BeleskaSlajda{09_1-s009}{9}{PUN i PDN mreže}{beleske/s009.tex}
\BeleskaSlajda{09_1-s010}{10}{PDN — moguće realizacije}{beleske/s010.tex}
\BeleskaSlajda{09_1-s011}{11}{PUN — moguće realizacije}{beleske/s011.tex}
\BeleskaSlajda{09_1-s012}{12}{Invertor sa pasivnim opterećenjem}{beleske/s012.tex}
\BeleskaSlajda{09_1-s013}{13}{Implicitna karakteristika prenosa}{beleske/s013.tex}
\BeleskaSlajda{09_1-s014}{14}{Granica zasićenja pogonskog tranzistora}{beleske/s014.tex}
\BeleskaSlajda{09_1-s015}{15}{Oblasti rada i pojačanje}{beleske/s015.tex}
\BeleskaSlajda{09_1-s016}{16}{Izlazni naponski nivoi}{beleske/s016.tex}
\BeleskaSlajda{09_1-s017}{17}{Aproksimacija napona logičke nule}{beleske/s017.tex}
\BeleskaSlajda{09_1-s018}{18}{Određivanje ulaznih pragova}{beleske/s018.tex}
\BeleskaSlajda{09_1-s019}{19}{Prag VIL: implicitno diferenciranje}{beleske/s019.tex}
\BeleskaSlajda{09_1-s020}{20}{Prag VIH: triodna oblast}{beleske/s020.tex}
\BeleskaSlajda{09_1-s021}{21}{Prag VIH: dve jednačine, dve nepoznate}{beleske/s021.tex}
\BeleskaSlajda{09_1-s022}{22}{Ispit: postupak analize}{beleske/s022.tex}
\BeleskaSlajda{09_1-s023}{23}{Održivost naponskih nivoa}{beleske/s023.tex}
\BeleskaSlajda{09_1-s024}{24}{Prag odlučivanja: kvadratna jednačina}{beleske/s024.tex}
\BeleskaSlajda{09_1-s025}{25}{Prag odlučivanja: približan račun}{beleske/s025.tex}
\BeleskaSlajda{09_1-s026}{26}{Ulazne struje i faktor grananja}{beleske/s026.tex}
\BeleskaSlajda{09_1-s027}{27}{Izlazni strujni kapaciteti}{beleske/s027.tex}
\BeleskaSlajda{09_1-s028}{28}{Dinamički režim: punjenje i pražnjenje}{beleske/s028.tex}
\BeleskaSlajda{09_1-s029}{29}{Uslov zasićenja pri pražnjenju}{beleske/s029.tex}
\BeleskaSlajda{09_1-s030}{30}{Dinamička otpornost MOS tranzistora}{beleske/s030.tex}
\BeleskaSlajda{09_1-s031}{31}{Srednja otpornost iz dve radne tačke}{beleske/s031.tex}
\BeleskaSlajda{09_1-s032}{32}{Aproksimacija srednje otpornosti}{beleske/s032.tex}
\BeleskaSlajda{09_1-s033}{33}{Integralna srednja otpornost}{beleske/s033.tex}
\BeleskaSlajda{09_1-s034}{34}{Ekvivalentna otpornost PMOS tranzistora}{beleske/s034.tex}
\BeleskaSlajda{09_1-s035}{35}{Aktivno opterećenje: izbor parametara}{beleske/s035.tex}
\BeleskaSlajda{09_1-s036}{36}{NMOS opterećenje sa kontrolnim naponom}{beleske/s036.tex}
\BeleskaSlajda{09_1-s037}{37}{Režim rada tranzistora opterećenja}{beleske/s037.tex}
\BeleskaSlajda{09_1-s038}{38}{Aktivno omsko opterećenje: logička jedinica}{beleske/s038.tex}
\BeleskaSlajda{09_1-s039}{39}{Početak provođenja pogonskog tranzistora}{beleske/s039.tex}
\BeleskaSlajda{09_1-s040}{40}{Prag VIL sa aktivnim omskim opterećenjem}{beleske/s040.tex}
\BeleskaSlajda{09_1-s041}{41}{Poseban slučaj i kvadratna jednačina}{beleske/s041.tex}
\BeleskaSlajda{09_1-s042}{42}{Rešenje za prag VIL}{beleske/s042.tex}
\BeleskaSlajda{09_1-s043}{43}{Oba tranzistora u omskoj oblasti}{beleske/s043.tex}
\BeleskaSlajda{09_1-s044}{44}{Prag VIH sa omskim opterećenjem}{beleske/s044.tex}
\BeleskaSlajda{09_1-s045}{45}{Rezultati za prag VIH}{beleske/s045.tex}
\BeleskaSlajda{09_1-s046}{46}{Logička nula sa omskim opterećenjem}{beleske/s046.tex}
\BeleskaSlajda{09_1-s047}{47}{Odnos dimenzija tranzistora}{beleske/s047.tex}
\BeleskaSlajda{09_1-s048}{48}{Karakteristika prenosa i kompromisi}{beleske/s048.tex}
\BeleskaSlajda{09_1-s049}{49}{Gejt opterećenja na napajanju}{beleske/s049.tex}
\BeleskaSlajda{09_1-s050}{50}{Logička jedinica sa zasićenim opterećenjem}{beleske/s050.tex}
\BeleskaSlajda{09_1-s051}{51}{Oba tranzistora u zasićenju}{beleske/s051.tex}
\BeleskaSlajda{09_1-s052}{52}{Konstantno pojačanje u oblasti zasićenja}{beleske/s052.tex}
\BeleskaSlajda{09_1-s053}{53}{Prelazak u omsku oblast i prag VIH}{beleske/s053.tex}
\BeleskaSlajda{09_1-s054}{54}{Rezultati za zasićeno opterećenje}{beleske/s054.tex}
\BeleskaSlajda{09_1-s055}{55}{Logička nula sa zasićenim opterećenjem}{beleske/s055.tex}
\BeleskaSlajda{09_1-s056}{56}{Uticaj zajedničke osnove}{beleske/s056.tex}
\BeleskaSlajda{09_1-s057}{57}{Obuhvat analize aktivnog opterećenja}{beleske/s057.tex}
\end{document}
